% stephane [DOT] adjemian [AT] univ [DASH] lemans [DOT] fr
\documentclass[10pt,a4paper,notitlepage,twocolumn]{article}
\usepackage{amsmath}
\usepackage{amssymb}
\usepackage{amsbsy}
\usepackage{bm}
\usepackage[T1]{fontenc}
\usepackage[utf8x]{inputenc}
\usepackage{palatino}
\usepackage{scrtime}
\usepackage[french]{babel}
\usepackage{float}
\usepackage{nicefrac}

%%%%%%%%%%%%%%%%%%%%%%%%%%%%%%%%%%%%%%%%%%%%%%%%%%%%%%%%%%%%%%%%%%%%%%%%%%%%%%%%%%%%%%%%%%%%%%%%%%%%%

\newcommand{\exercice}[1]{\textsc{\textbf{Exercice}} #1}
\newcommand{\question}[1]{\textbf{(#1)}}
\setlength{\parindent}{0cm}



\begin{document}


\title{\textsc{Séries temporelles}}
\author{\textsc{Université du Mans (Examen, L3)}}
\date{}

\pagenumbering{gobble}

\maketitle

\exercice{1} Soit le processus AR(2) défini par~:
\[
  Y_t = \frac{3}{8} + \frac{3}{4}Y_{t-1} - \frac{1}{8}Y_{t-2} + \varepsilon_t
\]
avec $\varepsilon_t$ un bruit blanc d'espérance nulle et de variance
$\sigma_{\varepsilon}^2$.\newline

\question{1} Montrer que ce processus est asymptotiquement stationnaire
au second ordre.\newline

\question{2} Calculer son espérance.\newline

\question{3} Écrire les équations de Yule-Walker et en déduire les
expressions des autocovariances $\gamma(0)$, $\gamma(1)$ et
$\gamma(2)$ en fonction de $\sigma_{\varepsilon}^2$.\newline

\question{4} Donner la relation de récurrence vérifiée par
l'autocovariance $\gamma(h)$ pour tout $h\geq 1$, puis en déduire la
forme générale de $\gamma(h)$.\newline

\bigskip

\exercice{2} Supposons que $\{y_t,t\in\mathbb Z\}$ soit un ARMA($1,1$)
de la forme~:
\[
y_t = 1 - \frac{1}{2}y_{t-1} + \varepsilon_t + \frac{1}{4} \varepsilon_{t-1}
\]
avec $\varepsilon_t$ un bruit blanc d'espérance nulle et de variance
1.\newline

\question{1} Le processus est-il asymptotiquement stationnaire au
second ordre et inversible~? Justifiez votre réponse.\newline

On suppose que les conditions initiales sont telles que le processus
est stationnaire au second ordre.\newline

\question{2} Quelles sont les implications de cette hypothèse sur les
moments d'ordre 1 et 2~? \question{3} Calculez l'espérance (on notera
$\mu$). \question{4} Calculez les autocovariances $\gamma(0)$ et
$\gamma(1)$. \question{5} Établir la relation de récurrence vérifiée
par $\gamma(h)$ pour $h\geq 2$ et en déduire une expression
explicite. \question{6} Que peut-on dire du signe de la fonction
d'autocorrélation $\rho(h)$ lorsque $h$ varie~?\newline

\bigskip

\exercice{3} Soit le processus AR(2) gaussien~:
\[
y_t = c + \varphi_1 y_{t-1} + \varphi_2 y_{t-2} + \varepsilon_t
\]
avec $\varepsilon_t \sim \mathcal{N}(0,\sigma_{\varepsilon}^2)$,
indépendants entre eux. On note $\mathcal Y_T \equiv
\{y_1,y_2,\dots,y_T\}$ l'échantillon observé et
$\bm{\theta}=(c,\varphi_1,\varphi_2,\sigma_{\varepsilon}^2)'$ le
vecteur des paramètres. On supposera dans tout l'exercice que le
processus est stationnaire.\newline

\question{1} En utilisant le théorème de Bayes, écrire la
décomposition de la vraisemblance exacte. Pourquoi cette vraisemblance
est-elle plus délicate à manipuler que dans le cas d'un AR(1)~?\newline

\question{2} Définir et écrire la vraisemblance conditionnelle (on
conditionnera par rapport aux deux premières observations $y_1$ et
$y_2$). En quoi diffère-t-elle de la vraisemblance exacte et quel est
son intérêt pratique~?\newline

\question{3} Donner l'expression de la log-vraisemblance
conditionnelle. Montrer qu'à $\sigma_{\varepsilon}^2$ fixé, la
maximisation de la log-vraisemblance conditionnelle par rapport à
$(c,\varphi_1,\varphi_2)$ équivaut à un problème de moindres carrés
ordinaires. En déduire l'estimateur correspondant sous forme
matricielle.\newline

\bigskip

\exercice{4} Soit $\{x_t,t\in\mathbb Z\}$ un processus AR(1)
stationnaire~:
\[
x_t = \rho x_{t-1} + \eta_t,\qquad |\rho|<1
\]
où $\eta_t$ est un bruit blanc d'espérance nulle et de variance
$\sigma_{\eta}^2$. On observe en réalité $x_t$ avec une erreur de
mesure~:
\[
y_t = x_t + \varepsilon_t
\]
où $\varepsilon_t$ est un bruit blanc d'espérance nulle et de variance
$\sigma_{\varepsilon}^2$, indépendant du processus
$\{\eta_s,\,s\in\mathbb Z\}$.\newline

\question{1} En utilisant le polynôme retard, montrer que le processus
$y_t$ vérifie~:
\[
(1-\rho L) y_t = \eta_t + \varepsilon_t - \rho\varepsilon_{t-1}
\]

\question{2} Calculer la fonction d'autocovariance du membre de droite
$\xi_t \equiv \eta_t + \varepsilon_t - \rho\varepsilon_{t-1}$. En
déduire que $\xi_t$ admet une représentation MA(1) du type $\xi_t =
u_t - \theta u_{t-1}$ où $u_t$ est un bruit blanc de variance
$\sigma_u^2$.\newline

\question{3} En déduire que $y_t$ est un processus ARMA(1,1) que l'on
explicitera. Donner le système d'équations vérifié par les paramètres
$\theta$ et $\sigma_u^2$ de la représentation.\newline

\question{4} Le processus $y_t$ est-il stationnaire au second ordre~?
Est-il inversible~? Justifier.

\end{document}

%%% Local Variables:
%%% mode: latex
%%% TeX-master: t
%%% End:
